\documentclass[a4paper,10pt]{article}

\usepackage[T1]{fontenc}
\usepackage[utf8]{inputenc}
\usepackage[ngerman]{babel}
\usepackage[top=1.5cm, bottom=1.5cm, left=2.3cm, right=2.3cm]{geometry}
\usepackage{listings}
\usepackage{xcolor}
\usepackage{parskip}
\usepackage{microtype}
\usepackage{booktabs}
\usepackage{helvet}     % Helvetica-compatible sans (~ Liberation Sans)
\usepackage{courier}    % monospace (~ DejaVu Sans Mono)
\renewcommand{\familydefault}{\sfdefault}

% Colors
\definecolor{codebg}{RGB}{244,248,240}
\definecolor{codegreen}{RGB}{26,58,26}
\definecolor{headcolor}{RGB}{46,64,87}
\definecolor{notecolor}{RGB}{240,247,255}
\definecolor{noteborder}{RGB}{100,149,200}

% Code style
\lstdefinelanguage{Cstar}{
  keywords={auto, if, else, while, return, void, uint64_t, char},
  keywordstyle=\bfseries\color{headcolor},
  comment=[l]{//},
  commentstyle=\itshape\color{gray},
  stringstyle=\color{codegreen},
  morestring=[b]",
  sensitive=true
}

\lstset{
  language=Cstar,
  backgroundcolor=\color{codebg},
  basicstyle=\ttfamily\footnotesize,
  frame=single,
  framerule=0.4pt,
  rulecolor=\color{lightgray},
  xleftmargin=12pt,
  xrightmargin=12pt,
  framexleftmargin=8pt,
  framexrightmargin=8pt,
  framextopmargin=5pt,
  framexbottommargin=5pt,
  breaklines=true,
  showstringspaces=false,
  tabsize=4,
  aboveskip=6pt,
  belowskip=6pt,
}

\setlength{\parskip}{2pt}
\setlength{\parindent}{0pt}

\begin{document}

% Header
{\small [Florian Schroffner] \hfill \today}

\vspace{0.2cm}
\hrule height 1.5pt
\vspace{0.15cm}

{\large \textbf{\textcolor{headcolor}{Projektvorschlag: Selfie/C* um Type Inference (\texttt{auto}) erweitern}}}\\

\vspace{0.15cm}
\hrule height 0.4pt
\vspace{0.2cm}

\textbf{Typsystem}

Die Erweiterung bleibt bewusst innerhalb des Typsystems von C* (\texttt{uint64\_t}, \texttt{uint64\_t*}): \texttt{auto} leitet einen dieser beiden bestehenden Typen ab, ohne neue Typen oder implizite Konversionen einzuführen. Der interessante Inferenzfall liegt nicht bei primitiven Typen, sondern bei \textbf{Pointern}: \texttt{\&}, \texttt{*} und Pointer-Arithmetik verändern den Typ --- und genau diesen Typ liefert Selfies Ausdrucks-Compiler bereits als Rückgabewert zurück. Vergleiche liefern \texttt{uint64\_t} mit Wert 0 oder 1, analog zu C:

\vspace{1pt}
\begin{center}
\small
\begin{tabular}{lll}
\toprule
\textbf{Operation} & \textbf{Inferenz-Beispiel} & \textbf{Ergebnis} \\
\midrule
Integer-Literal & \texttt{auto n = 10;}      & \texttt{uint64\_t}  \\
Adresse \texttt{\&} & \texttt{auto p = \&n;}  & \texttt{uint64\_t*} \\
String-Literal  & \texttt{auto s = "Selfie";} & \texttt{uint64\_t*} \\
Deref \texttt{*} & \texttt{auto x = *p;}     & \texttt{uint64\_t}  \\
Ptr-Arithmetik  & \texttt{auto q = p + 1;}   & \texttt{uint64\_t*} \\
Vergleich       & \texttt{auto b = (n > 0);} & \texttt{uint64\_t} (0/1) \\
Fehler          & \texttt{p + q} (beide Ptr) & \textit{Compilerfehler} \\
\bottomrule
\end{tabular}
\end{center}
\vspace{1pt}

\textbf{Beispielprogramm}

\begin{lstlisting}
uint64_t increment(uint64_t i) { return i + 1; }

uint64_t main() {
    auto n      = 10;             // uint64_t  -- Integer-Literal
    auto result = increment(n);   // uint64_t  -- Rueckgabetyp der Funktion
    auto flag   = (result > 0);   // uint64_t  -- Vergleich: 0 oder 1

    auto p      = &n;             // uint64_t* -- Adresse von n
    auto q      = p + 1;          // uint64_t* -- Pointer-Arithmetik
    auto x      = *p;             // uint64_t  -- Dereferenz

    auto s      = "Selfie";       // uint64_t* -- String-Literal (char*)
    // auto err = p + q;          // FEHLER    -- Pointer + Pointer ist undefiniert
    return x - n;                 // 0 bei Erfolg
}
\end{lstlisting}

\textbf{Architektur: weiterhin ein einziger Pass}

Die Erweiterung verändert Selfies Single-Pass-Architektur \emph{nicht}: \texttt{Scanner $\rightarrow$ rekursiver Abstiegsparser \texttt{+} Code-Emitter $\rightarrow$ RISC-U}. Es wird \textbf{kein} expliziter AST aufgebaut und \textbf{kein} eigener Typcheck-Pass eingeführt. Der Schlüssel ist, dass Selfies Ausdrucks-Compiler bereits lokal Typen verfolgt: Jede \texttt{compile\_*}-Funktion gibt den Typ (\texttt{UINT64\_T} oder \texttt{UINT64STAR\_T}) des soeben erzeugten Werts zurück. Eine \texttt{auto}-Deklaration \emph{liest diesen Typ lediglich aus} dem Rückgabewert von \texttt{compile\_expression()} und trägt ihn in die Symboltabelle ein. Die Änderung betrifft drei Funktionen in \texttt{selfie.c}; der Code-Generator und das Typsystem bleiben unberührt.

\textbf{Grenzen der Inferenz}

Die Inferenz ist bewusst \textbf{lokal} gehalten: Der Typ einer \texttt{auto}-Variable wird ausschließlich aus dem direkten Initialisierer abgeleitet --- nie aus späterem Verwendungskontext. Da im Single-Pass ein Bezeichner vor seiner Verwendung deklariert sein muss, sind zirkuläre und vorwärtsgerichtete Abhängigkeiten (\texttt{auto x = y; auto y = x;}) automatisch ein Fehler (``undeclared identifier''). Damit bleibt die Inferenz stets entscheidbar, auf Kosten von Ausdrucksstärke.

\end{document}
